% =============================================================================
% CHAPTER 9: ADAPTIVE DIGITAL TWIN
% IP Traceability: IP-008 (digitaltwin.pdf)
% =============================================================================

\chapter{Adaptive Digital Twin}
\label{chap:digital-twin}

\section{IP Document Summary}

\begin{tracerecord}
\textbf{IP Source ID} & \ipid{IP-008} \\
\textbf{Document} & digitaltwin.pdf \\
\textbf{Author} & Tom Pfaff, Craig Turrell \\
\textbf{Date} & January 2026 \\
\textbf{Pages} & 11 \\
\textbf{Abstract} & Adaptive world models for real-world systems using Liquid Neural Networks \\
\end{tracerecord}

\subsection{Document Abstract}

\begin{ipsourcebox}{IP-008, Abstract}
\textit{``We present an adaptive world model architecture that enables autonomous systems to build, maintain, and reason with internal representations of physical reality. Unlike static simulation environments or pre-trained models, our framework creates a living digital twin---a continuously-updated internal model that learns the dynamics of the real world through embodied experience, adapts to changing conditions, and enables predictive reasoning about physical interactions.''}
\end{ipsourcebox}

\section{Key Concepts Identified}

\begin{table}[H]
\centering
\caption{IP-008 Key Concepts}
\begin{tabularx}{\textwidth}{@{}p{3.5cm}Xl@{}}
\toprule
\textbf{Concept} & \textbf{Description} & \textbf{Section} \\
\midrule
Liquid Neural Networks & Continuous-time adaptive dynamics & \ref{sec:ip008-lnn} \\
Physics World Model & Learned dynamics for prediction & \ref{sec:ip008-physics} \\
Lyapunov Stability & Mathematical stability guarantees & \ref{sec:ip008-lyapunov} \\
Adaptive Learning & Online model updates & \ref{sec:ip008-adaptive} \\
\bottomrule
\end{tabularx}
\end{table}

\section{Concept: Liquid Neural Networks}
\label{sec:ip008-lnn}

\begin{ipsourcebox}{IP-008, Liquid Neural Networks Section}
\textit{``Liquid Time-Constant Cell implements continuous-time dynamics:
$$\frac{dx}{dt} = \frac{-x + f(W_{in} \cdot input + W_{rec} \cdot x + b)}{\tau(x, input)}$$
Where $\tau$ is a state-dependent time constant allowing neurons to respond quickly to important signals (low $\tau$) while maintaining stable states under other conditions (high $\tau$).''}
\end{ipsourcebox}

\begin{tracerecord}
\textbf{IP Source ID} & IP-008 \\
\textbf{IP Location} & Liquid Neural Networks Section \\
\textbf{Concept} & Continuous-time dynamics with adaptive time constants \\
\midrule
\textbf{Implementation Path} & \filepath{src/intelligence/analytics/} \\
\textbf{Adaptation Type} & Conceptual (adaptive behavior patterns) \\
\end{tracerecord}

\begin{adaptbox}
The IP paper describes Liquid Neural Networks with adaptive time constants. While SAP-OSS does not implement LNNs directly, the concept of adaptive response---responding differently to significant vs. routine inputs---is captured in \funcname{\_determine\_method()}, which adapts the detection method based on data characteristics (normal vs. non-normal distributions).
\end{adaptbox}

\section{Concept: Physics World Model}
\label{sec:ip008-physics}

\begin{ipsourcebox}{IP-008, Adaptive Digital Twin Section}
\textit{``The physics world model is the core of our adaptive digital twin---an internal representation of physical reality that learns, adapts, and enables reasoning. The PhysicsDynamicsModel learns $p(s_{t+1} | s_t, a_t)$ using residual prediction (predicts $\Delta s$, not $s'$) with uncertainty estimation.''}
\end{ipsourcebox}

\begin{tracerecord}
\textbf{IP Source ID} & IP-008 \\
\textbf{IP Location} & Physics World Model Section \\
\textbf{Concept} & Learned predictive model with uncertainty \\
\midrule
\textbf{Implementation Path} & \filepath{src/intelligence/analytics/anomaly\_detector.py} \\
\textbf{Adaptation Type} & Conceptual (predictive baseline for anomaly detection) \\
\end{tracerecord}

\begin{adaptbox}
The IP paper describes a physics world model that learns expected behavior and detects deviations. This concept is adapted to financial analytics: the \funcname{AnomalyDetector} builds a statistical model of expected values (mean, variance) and detects deviations (anomalies) from this baseline. The principle of ``learning what's normal to detect what's abnormal'' is preserved.
\end{adaptbox}

\section{Concept: Lyapunov Stability}
\label{sec:ip008-lyapunov}

\begin{ipsourcebox}{IP-008, Stability-Aware Control Section}
\textit{``The Lyapunov controller provides stability guarantees by learning both a policy $\pi(s)$ and a Lyapunov function $V(s)$ such that:
\begin{enumerate}
\item $V(s) > 0$ for all $s \neq 0$ (positive definite)
\item $dV/dt < 0$ along trajectories (decreasing)
\end{enumerate}
This ensures the system converges to a stable equilibrium.''}
\end{ipsourcebox}

\begin{tracerecord}
\textbf{IP Source ID} & IP-008 \\
\textbf{IP Location} & Stability-Aware Control Section \\
\textbf{Concept} & Mathematical stability guarantees \\
\midrule
\textbf{Implementation Path} & Conceptual influence \\
\textbf{Adaptation Type} & Design principle (bounded behavior) \\
\end{tracerecord}

\begin{adaptbox}
The IP paper emphasizes Lyapunov stability for bounded, predictable behavior. While SAP-OSS does not implement Lyapunov controllers, the principle of bounded behavior is captured in the anomaly detection thresholds: scores are bounded, classifications are finite (spike/drop/statistical/none), and detection methods are well-defined. This provides predictable, bounded output from the analytics system.
\end{adaptbox}

\section{Concept: Adaptive Learning}
\label{sec:ip008-adaptive}

\begin{ipsourcebox}{IP-008, Living Digital Twin Section}
\textit{``Unlike static digital twins that require expert construction, our adaptive digital twin learns from the robot's own experiences:
\begin{itemize}
\item Starts with minimal assumptions
\item Improves with every interaction
\item Adapts when the environment changes
\item Captures the specific dynamics of \textit{this} robot in \textit{this} environment
\end{itemize}''}
\end{ipsourcebox}

\begin{tracerecord}
\textbf{IP Source ID} & IP-008 \\
\textbf{IP Location} & Living Digital Twin Section \\
\textbf{Concept} & Online learning from experience \\
\midrule
\textbf{Implementation Path} & \filepath{src/intelligence/analytics/anomaly\_detector.py} \\
\textbf{Adaptation Type} & Conceptual (per-batch adaptation) \\
\end{tracerecord}

\begin{adaptbox}
The IP paper describes continuous learning from experience. The SAP-OSS \funcname{AnomalyDetector} adapts this concept through per-batch analysis: each batch of data triggers a normality test, and the detection method adapts to the current data characteristics. While not full online learning, this captures the principle of responding to the specific data being analyzed rather than using fixed parameters.
\end{adaptbox}

\section{Implementation Summary}

\begin{table}[H]
\centering
\caption{IP-008 Concept Implementation Summary}
\begin{tabularx}{\textwidth}{@{}p{3cm}p{3.5cm}Xc@{}}
\toprule
\textbf{Concept} & \textbf{IP Reference} & \textbf{Implementation} & \textbf{Status} \\
\midrule
Liquid Neural Networks & LNN Section & Adaptive method selection & Yes \\
Physics World Model & World Model Section & Statistical baseline modeling & Yes \\
Lyapunov Stability & Control Section & Bounded output design & Yes \\
Adaptive Learning & Living Twin Section & Per-batch adaptation & Yes \\
\bottomrule
\end{tabularx}
\end{table}

\paragraph{Note:}
IP-008 describes robotics and physical systems. The SAP-OSS adaptation applies the underlying principles (adaptive response, learned baselines, bounded behavior, continuous improvement) to financial analytics rather than physical control.